\section{\label{II_1_B} Outils de segmentation automatique}

Différents projets, extérieurs au domaine du patrimoine, tentent de mettre en place des algorithmes de segmentation vidéo. Bien que ces projets mettent en place l'analyse automatique de corpus vidéos moins spécifiques que ceux constitués d'\gls{Archivesaudiovisuelles}, leur architecture donne des pistes d'expérimentations afin de pouvoir implémenter un tel processus automatique au sein d'un logiciel de gestion des collections. En effet, face aux limites des méthodes d'\gls{Apprentissage profond}, nous tenterons de montrer dans quelle mesure la \gls{Vision par ordinateur} classique, fondée sur une architecture plus stable, peut s'appliquer à la tâche automatique que nous souhaitons implémenter, la segmentation d'un flux vidéo en différentes séquences. L'optique de cette étude est de fournir une proposition d'implémentation de fonctionnalités automatiques au sein d'un outil de gestion des collections d'un point de vue métier, en prenant en compte les besoins utilisateurs.

\subsection{Décomposer le mouvement}

Tout d'abord, afin de pouvoir implémenter des algorithmes plus légers que ceux des \gls{Transformers}, une des solutions envisagée consisterait à transformer un flux d'images en mouvement en la série d'images fixes qui le composent. Ce processus écarterait notamment le besoin d'analyser les relations entre les images par un mécanisme d'auto-attention prévu par les \gls{Transformers}. En effet, le médium cinématographique se caractérise par l'enchaînement d'une succession d'images fixes produisant l'illusion du mouvement. La transformation de ce flux à l'état d'images fixes par un outil technique permettrait de pouvoir gérer l'analyse sémantique de ces images par des processus légers prévus par la \gls{Vision par ordinateur} classique. Ces images fixes pourraient ensuite être regroupées en séquences après une analyse de leurs différences syntaxiques. Grâce à ce procédé, pourraient être appliqués les mêmes procédés d'analyse que ceux prévus pour l'image fixe, fondés sur la \gls{Vision par ordinateur} classique. 

Cette segmentation d'un flux audiovisuel fondé sur les différences de syntaxe entre les \gls{plan} serait la condition préalable d'une analyse sémantique, regroupant ces \gls{plan} fixes en des \gls{sequence} sémantiques. Dans un article sur la mise en place d'un processus de \gls{segtemp}, Yuelin Xin prône notamment le découpage syntaxique tel que nous l'avons défini comme étape préliminaire à l'analyse sémantique des images par des modèles d'intelligence artificielle \footcite{xin2023}. La segmentation syntaxique d'un flux audiovisuel en différents \gls{plan} fixes servirait alors de filtre permettant la classification sémantique de ces \gls{plan}. 

Ainsi, la segmentation automatique d'un document audiovisuel en \gls{sequence}s serait fondée sur l'analyse préalable des différences entre les \gls{plan} fixes qui la composent. Ainsi, si la syntaxe visuelle de deux \gls{plan} fixes est jugée suffisamment différente selon un seuil mathématique prédéfini, ces \gls{plan} appartiendront à des \gls{sequence} distinctes\footcite{xin2023}. Ces \gls{plan} sont en réalité retenus par l'algorithme comme des \gls{imgcle} contenues au sein de chaque \gls{sequence}, calculées comme l'image la plus représentative de la présence des éléments dans un flux donné \footcite{xin2023}. Ce principe, qualifié de sélection de données, permet d'alléger nettement la puissance de calcul et la structure des algorithmes \footcite{xin2023}. Ces \gls{imgcle} seraient ensuite analysées par un algorithme fondé sur des \gls{cnn}, afin d'y ajouter une analyse sémantique et de les regrouper en \gls{sequence}s. 

\subsection{PySceneDetect}

Ce processus est notamment appliqué par la bibliothèque python \gls{pyscene} qui permet la segmentation d'un flux vidéo à l'aide de différents paramètres définis en amont \footcite{zotero-item-706}. Les images extraites peuvent ensuite être regroupées de manière thématique, une fois l'analyse de leurs différences syntaxiques effectuée. Sans toutefois extraire la totalité des images fixes qui composent une \gls{sequence} audiovisuelle, \gls{pyscene} permet d'en extraire les \gls{imgcle}, c'est-à-dire les principaux différents \gls{plan}.

\gls{pyscene} est un outil en accès libre fondé sur des tâches de \gls{Vision par ordinateur} classique (non fondé sur un algorithme d'\gls{ia}), qui concrétise la mise en oeuvre de la comparaison syntaxique d'images fixes. L'outil s'appuie sur des détecteurs mathématiques afin d'analyser les différences structurelles entre les \gls{imgcle} tirées d'une vidéo \footcite{zotero-item-708}. Il s’agit ainsi d’un outil basé sur des seuils et mesures visuelles mathématiques paramétrables, et ne nécesite pas d'entraînement particulier, comme c'est le cas des outils de \gls{Vision par ordinateur} fondés sur l'\gls{Apprentissage profond}. \gls{pyscene} permet d'isoler des \gls{imgcle} de l'image animée, de les regrouper par similarité et de les dissocier selon leurs différences. 

\subsection{Fonctionnement}

\subsubsection{Paramétrages}
Trois méthodes de détection sont proposées par l’outil : une détection des fondus entrants et sortants (méthode \textit{treshold}, qui repère les niveaux de noir dans l’image, une détection de coupes rapides qui repère les changements de contenu entre images (méthode \textit{detect-content}), et une détection de coupes fondée sur la différence d’un plan avec les \gls{plan} adjacents (méthode \textit{detect-adaptative}) \footcite{zotero-item-708}. 
En théorie, \gls{pyscene} isole des \gls{plan} fixes, correspondant à des ruptures syntaxiques dans un flux d'images, et non des \gls{sequence}s, unités de sens narratives. Néanmoins, \gls{pyscene} permet de contourner cette limite grâce à la flexibilité de son paramétrage. 
En effet, l'élévation de la valeur du seuil de détection ou la fixation d'une durée minimale de \gls{sequence} (\textit{min-scene-len}), l'outil est moins sensibles aux variations minimales entre images et ne déclenche une rupture qu'en cas de déctection de changements visuels majeurs tels que l'éclairage, le décor, la composition, etc \footcite{zotero-item-708}. Par ce biais, la segmentation syntaxique proposée par \gls{pyscene} pourrait permettre de diviser un flux vidéo en différentes \gls{sequence}s, que la variation visuelle permettrait de regrouper de manière sémantique. Une fois le processus finalisé, l'outil génère des marqueurs temporels sur les images extraites et enregistre automatiquement \gls{imgcle} associée à chaque \gls{sequence} retenue. Le diagramme de flux \ref{fig:pyscenedetect_pipeline} ci-dessous illustre la manière dont l'outil pourrait être paramétré afin de segmenter le flux vidéo en \gls{sequence}s.

\begin{figure}[H]
	\centering
	\begin{tikzpicture}[node distance=1.5cm, auto, >=stealth,
		block/.style={rectangle, draw, fill=blue!5, text width=2.6cm, text centered, rounded corners, minimum height=1.2cm, font=\small},
		decision/.style={diamond, draw, fill=orange!10, text width=2cm, text centered, inner sep=0pt, font=\small},
		line/.style={draw, ->, thick}]
		
		\node [block] (video) {\textbf{Flux Vidéo Continu}\\};
		\node [block, right=1cm of video] (subss) {\textbf{Extraction de plans}\\};
		\node [block, right=1cm of subss] (diff) {\textbf{Calcul d'écart}};
		\node [decision, below=1cm of diff] (seuil) {\textbf{Écart supérieur au Seuil ?}};
		\node [block, left=1cm of seuil] (continue) {\textbf{Même séquence}\\};
		\node [block, right=1.2cm of seuil] (cut) {\textbf{Nouvelle séquence}\\};
		\node [block, below=0.8cm of cut] (imgcle) {\textbf{Sauvegarde d'images-clés et marqueurs temporels}};
		
		
		% Liens
		\path [line] (video) -- (subss);
		\path [line] (subss) -- (diff);
		\path [line] (diff) -- (seuil);
		\path [line] (seuil) -- node [above] {Non} (continue);
		\path [line] (seuil) -- node [above] {Oui} (cut);
		\path [line] (cut) -- (imgcle);
		
	\end{tikzpicture}
	\caption{Schéma du processus de détection de coupures par \gls{pyscene}}
	\label{fig:pyscenedetect_pipeline}
\end{figure}

Pour chaque changement détecté, une nouvelle \gls{sequence} est ainsi créée et une \gls{imgcle} est automatiquement enregistrée, associée au fichier vidéo d'origine. Nous proposerons la réutilisation de cette image au sein du processus suivant, dédié à l'analyse sémantique des séquences préalablement détectées, afin de ne pas surcharger les algorithmes. Ce système permettrait ainsi d’alléger le processus d’analyse sémantique, puisque les algorithmes reçoivent des images fixes au lieu de fichiers vidéos plus ou moins lourds. 

\subsubsection{Limites}

Toutefois, certaines mises en application de l'outil soulignent ses limites. Dans un article sur l'\gls{Apprentissage profond} comme outil d'analyse du cinéma des premiers temps, Samarth Bhargav explique par exemple que l'outil détecte beaucoup plus de coupures de \gls{plan}s que nécessaire au sein de ce type de corpus \footcite{bhargav2019}. Ce problème peut être lié à la qualité amoindrie de ce type d'archives, même après numérisation ou restauration \footcite{bhargav2019}. En effet, la qualité de l'image détermine nécessairement la performance de l'outil d'analyse utilisé et des résultats reçus \footcite{norindr2023}. De plus, Samarth Bhargav souligne la nécessité de connaître de manière précise les modalités de paramétrage de l'outil, afin de pouvoir l'adapter adéquatement à la nature du corpus analysé \footcite{bhargav2019}. Toutefois, un des avantages majeurs de l'outil \gls{pyscene} est sa disponibilité et la possibilité de le mettre en oeuvre facilement et gratuitement. De plus, c'est un outil fondé sur des tâches classiques de \gls{Vision par ordinateur} analytiques qui ne nécessitent pas d'entraînement particulier. Son paramétrage est également documenté et facilement modifiable, contrairement à ceux des réseaux de neurones. 
